% Generated from locale/tr/content/first-order-logic/natural-deduction/identity.tex; do not hand-edit.


\olfileid[tr]{fol}{ntd}{ide}

\olsection{\usetoken{S}{identity} İçeren \usetoken{P}{derivation}}

!!^{derivation}s, !!{identity} içeriyorsa ek çıkarım kuralları gerektirir.

\begin{defish}
\AxiomC{}
\RightLabel{\Intro{\eq}}
\UnaryInfC{$\eq[t][t]$}
\DisplayProof
\hfill
\begin{tabular}{r}
\AxiomC{$\eq[t_1][t_2]$}
\AxiomC{$!A(t_1)$}
\RightLabel{\Elim{\eq}}
\BinaryInfC{$!A(t_2)$}
\DisplayProof
\\[3ex]
\AxiomC{$\eq[t_1][t_2]$}
\AxiomC{$!A(t_2)$}
\RightLabel{\Elim{\eq}}
\BinaryInfC{$!A(t_1)$}
\DisplayProof
\end{tabular}
\end{defish}

Yukarıdaki kurallarda $t$, $t_1$ ve $t_2$ kapalı terimlerdir.
\Intro{\eq} kuralı sayesinde $\eq[t][t]$ biçimindeki herhangi bir özdeşlik
ifadesini hiçbir varsayım olmaksızın doğrudan !!{derive} mümkündür.

\begin{ex}
$s$ ve $t$ kapalı terimlerse $!A(s), \eq[s][t] \Proves !A(t)$ olur:
\begin{prooftree}
\AxiomC{$\eq[s][t]$}
\AxiomC{$!A(s)$}
\RightLabel{$\Elim{\eq}$}
\BinaryInfC{$!A(t)$}
\end{prooftree}
Bu, ``özdeşlerin birbirinin yerine konabilirliği ilkesi'' veya Leibniz Yasası
adıyla biliniyor olabilir.
\end{ex}

\begin{prob}
$=$ bağıntısının hem simetrik hem de geçişli olduğunu ispatlayın; yani şu
iki ifade için !!{derivation}s verin:
$\lforall[x][\lforall[y][(\eq[x][y] \lif \eq[y][x])]]$ ve
$\lforall[x][\lforall[y][\lforall[z]((\eq[x][y] \land \eq[y][z]) \lif
\eq[x][z])]]$.
\end{prob}

\begin{ex}
Aşağıda !!{derive} eylemini gerçekleştirerek şu !!{sentence} ifadesine ulaşırız:
\begin{align*}
& \lforall[x][\lforall[y][((!A(x) \land !A(y)) \lif \eq[x][y])]]
\intertext{şu !!{sentence} ifadesinden:}
& \lexists[x][\lforall[y][(!A(y) \lif \eq[y][x])]].
\end{align*}
!!{derivation} üzerinde geriye doğru ilerleriz:
\begin{prooftree}
\AxiomC{$\lexists[x][\lforall[y][(!A(y) \lif \eq[y][x])]]
\quad \Discharge{!A(a) \land !A(b)}{1}$}
\DeduceC{$\eq[a][b]$}
\DischargeRule{\Intro{\lif}}{1}
\UnaryInfC{$((!A(a) \land !A(b)) \lif \eq[a][b])$}
\RightLabel{\Intro{\lforall}}
\UnaryInfC{$\lforall[y][((!A(a) \land !A(y)) \lif \eq[a][y])]$}
\RightLabel{\Intro{\lforall}}
\UnaryInfC{$\lforall[x][\lforall[y][((!A(x) \land !A(y)) \lif \eq[x][y])]]$}
\end{prooftree}
Şimdi ana varsayımı kullanmamız gerekir: bu bir varlıksal !!{formula}
olduğundan ara sonuç $\eq[a][b]$'yi !!{derive} üzere
\Elim{\lexists} kullanırız.
\begin{prooftree}
\AxiomC{$\lexists[x][\lforall[y][(!A(y) \lif \eq[y][x])]]$}
\AxiomC{$\Discharge{\lforall[y][(!A(y) \lif \eq[y][c]])}{2}$}
\noLine
\UnaryInfC{$\Discharge{!A(a) \land !A(b)}{1}$}
\DeduceC{$\eq[a][b]$}
\DischargeRule{\Elim{\lexists}}{2}
\BinaryInfC{$\eq[a][b]$}
\DischargeRule{\Intro{\lif}}{1}
\UnaryInfC{$((!A(a) \land !A(b)) \lif \eq[a][b])$}
\RightLabel{\Intro{\lforall}}
\UnaryInfC{$\lforall[y][((!A(a) \land !A(y)) \lif \eq[a][y])]$}
\RightLabel{\Intro{\lforall}}
\UnaryInfC{$\lforall[x][\lforall[y][((!A(x) \land !A(y)) \lif \eq[x][y])]]$}
\end{prooftree}
Sağ üstteki alt-!!{derivation}, varsayımlar kullanılıp $\eq[a][c]$ ile
$\eq[b][c]$ gösterilerek tamamlanır. Bunun için iki ayrı !!{derivation}s
gerekir. $\eq[a][c]$ için !!{derivation} şöyledir:
\begin{prooftree}
\AxiomC{$\Discharge{\lforall[y][(!A(y) \lif \eq[y][c]])}{2}$}
\RightLabel{\Elim{\lforall}}
\UnaryInfC{$!A(a) \lif \eq[a][c]$}
\AxiomC{$\Discharge{!A(a) \land !A(b)}{1}$}
\RightLabel{\Elim{\land}}
\UnaryInfC{$!A(a)$}
\RightLabel{\Elim{\lif}}
\BinaryInfC{$\eq[a][c]$}
\end{prooftree}
$\eq[a][c]$ ve $\eq[b][c]$ ifadelerinden \Elim{\eq} yoluyla
$\eq[a][b]$'yi !!{derive} mümkündür.
\end{ex}

\begin{prob}
Aşağıdaki !!{derivation}s verin; bunların hedefleri şu !!{formula}s olsun:
\begin{enumerate}
\item $\lforall[x][\lforall[y][((\eq[x][y] \land !A(x)) \lif !A(y))]]$
\item $\lexists[x][!A(x)] \land \lforall[y][\lforall[z][((!A(y) \land
    !A(z)) \lif \eq[y][z])]] \lif \lexists[x][(!A(x) \land
  \lforall[y][(!A(y) \lif \eq[y][x])])]$
\end{enumerate}
\end{prob}

